\section{Literature Review}

\subsection{Calculation of WACC}

For the WACC data, this study mainly uses the dataset from Calcaterra et al. \cite{calcaterra2024reducing}, which contains predictions for WACC across different regions. The calculation begins with analyzing past financial data at the enterprise level, calculating the cost of debt and equity, and incorporating the leverage ratio and tax rate to determine each company's WACC as follows \cite{calcaterra2024reducing}:
\begin{equation}
WACC_{it}=L_{it}\times r_{Dit}\times (1-\mathrm{TaxRate}_{it})+(1-L_{it})\times r_{Eit}.
\end{equation}
where $WACC_{it}$ represents the weighted average cost of capital for firm $i$ at time $t$, $L_{it}$ is the leverage ratio, $r_{Dit}$ is the cost of debt, $r_{Eit}$ is the cost of equity, and $\mathrm{TaxRate}_{it}$ is the tax rate for firm $i$ at time $t$.

The national WACC is then derived by weighting and averaging each company's WACC based on its share of the national utilities-market revenue \cite{calcaterra2024reducing}:
\begin{equation}
WACC_{ct}=\sum_i\left(\frac{\mathrm{Revenue}_{it}}{\mathrm{TotalRevenue}_{ct}}\right)\times WACC_{it}.
\end{equation}
where $WACC_{ct}$ represents the weighted average cost of capital for country $c$ at time $t$, $\mathrm{Revenue}_{it}$ is the revenue of firm $i$ at time $t$, and $\mathrm{TotalRevenue}_{ct}$ is the total utilities-market revenue of country $c$ at time $t$.

After that, Calcaterra et al. \cite{calcaterra2024reducing} introduce a time dimension and financial learning, simulating how technological progress and experience accumulation can lower WACC by reducing technological risks and the required safety margin. The financial-learning scenario models capital costs as a function of cumulative deployment:
\begin{equation}
\mathrm{CoC}_{tnT}=\mathrm{CoC}_{0n,T}\times\left(\frac{Y_{tn,T}}{Y_{0n,T}}\right)^{b_T}.
\end{equation}
where $\mathrm{CoC}_{tnT}$ is the cost of capital for country $n$ at time $t$ and technology $T$, $\mathrm{CoC}_{0n,T}$ is the initial cost of capital, $Y_{tn,T}$ is cumulative technology deployment, $Y_{0n,T}$ is initial technology deployment, and $b_T=\log_2(1-0.05)<0$ is the financial-learning exponent. The separate CoC-convergence scenario models convergence towards developed-country financing conditions by 2050.

Utilizing the aforementioned data and assumptions, it can employ a suite of climate-energy-economy models to simulate the WACC across various regions and energy sectors.

\subsection{Impact of Fossil Fuel Prices on WACC Components}

The previous subsection outlines the calculation of $WACC_{it}$ using firm-level data, in which the cost of equity ($r_{Eit}$) and the cost of debt ($r_{Dit}$) are not static. They are profoundly influenced by broader market dynamics, among which fossil fuel prices play a critical role \cite{melodia2022energy}.

The cost of equity ($r_{Eit}$) is estimated using the Capital Asset Pricing Model (CAPM), defining the expected return as:
\begin{equation}
r_{Eit}=r_{f,t}+\beta_i(r_{m,t}-r_{f,t}).
\end{equation}
where $r_{f,t}$ is the risk-free rate at time $t$, $\beta_i$ is the firm's systematic risk, and $(r_{m,t}-r_{f,t})$ is the equity risk premium (ERP).

Fossil fuel prices could influence the components of the CAPM equation. The substitution effect suggests that higher fossil fuel prices can enhance the economic competitiveness of renewable energy projects and influence perceived risk and the required return, $r_{Eit}$ \cite{Lieberman2007RenewableHedge}. Additionally, the impact on systematic risk ($\beta_i$) is ambiguous. Renewable assets can act as a hedge against fossil fuel price volatility due to their zero fuel cost \cite{Lieberman2007RenewableHedge}. Stock-market evidence varies: Vr\^inceanu et al. find no significant effect of oil-price volatility on renewable-stock return volatility in their sample \cite{vrinceanu2020influence}. Renewable and fossil power portfolios can also differ in financial risk and return, as reported for India \cite{su13052573}.

Similarly, the cost of debt ($r_{Dit}$) is also susceptible to fossil fuel price shocks, which can contribute to inflation \cite{melodia2022energy,millischer2024renewables}. Changes in the interest-rate environment affect economy-wide borrowing costs and WACC \cite{dobrowolski2022weighted}. At the project level, Power Purchase Agreements (PPAs), counterparty risk, and risk-transfer instruments can affect financing conditions and the credit spread component of $r_{Dit}$ \cite{duma2023risk}. Evidence on environmental financial premia is conditional: Wilberg et al. find no aggregate ESG-score effect but an association between the environmental component and lower capital costs \cite{kjellevoll2023impact}.

Therefore, the final calculated $WACC_{it}$ for a renewable energy firm is a function not only of its internal financial structure ($L_{it}$) but also of these external pressures that alter the fundamental costs of its capital ($r_{Eit}$ and $r_{Dit}$). The interplay of these mechanisms is summarized in Table~\ref{tab:channels}. Ultimately, the net impact on the WACC of renewable energy projects hinges on the relative magnitude of these opposing forces: the project-level advantages stemming from higher fossil fuel prices versus the detrimental consequences of macroeconomic contagion.

\begin{table}[t]
\centering
\caption{Theoretical Channels of Fossil Fuel Price Transmission to Renewable Energy WACC}
\label{tab:channels}
\scriptsize
\begin{tabular}{P{0.27\linewidth}P{0.27\linewidth}P{0.34\linewidth}}
\toprule
\textbf{Price Signal} & \textbf{Affected Component} & \textbf{Transmission Channel} \\
\midrule
Level increase & $r_E$ (cost of equity) & Substitution effect \\
Level/volatility & $r_E$ (cost of equity) & Market correlation \\
Volatility increase & $r_E$ (cost of equity) & Hedging property \\
Volatility increase & $r_E$ (cost of equity) & Climate risk pricing \\
Volatility increase & $r_D$ (cost of debt) & Macroeconomic contagion \\
Level increase & $r_D$ (cost of debt) & Off-taker stability \\
Level increase & $r_D$ (cost of debt) & PPA contract value \\
Volatility increase & $r_D$ (cost of debt) & Debt market green premium \\
\bottomrule
\end{tabular}
\end{table}

\subsection{AI-Based Forecasting for Energy and Climate Applications}

Recent ACAIT proceedings show a growing interest in using deep learning for environmental, energy, and carbon-related forecasting tasks. Xie et al. \cite{acaitXie2024carbon} study carbon-emission prediction using deep time-series models. Chen et al. \cite{acaitChen2023transfer} propose a transfer-based adaptive recurrent neural network for enterprise carbon-emission prediction.

These studies are relevant because they position AI forecasting as a decision-support tool for climate and energy management rather than as a purely technical prediction exercise. However, the two cited ACAIT studies focus on carbon emissions as the final prediction target. They do not explicitly connect fossil fuel price forecasts to macroeconomic transmission channels, regional financing conditions, or project-level WACC. This gap motivates the present study's three-stage design: fossil fuel prices are first predicted with a hybrid sequential model, inflation is then modeled as a macroeconomic covariate, and the resulting signals are finally linked to capital-cost sensitivity for energy transition projects.
